\section{Montrer l'application}

\begin{frame}{KYST : le flux tiré jusqu'à la réservation}
  \small
  \begin{columns}[T,onlytextwidth]
    \begin{column}{0.58\textwidth}
      {\Large\bfseries\color{mcmPrimary} kyst.tech}\\[0.15cm]
      \textbf{Keep Your Surgeries Timely} : l'application web qui applique le
      flux tiré en consultation, du planning statique à l'adaptation dynamique.

      \vspace{0.12cm}
      \begin{exampleblock}{Trois gestes}
        \begin{itemize}
          \item saisir une demande en consultation ;
          \item obtenir deux dates, justifiées ;
          \item confirmer : vacation et lit sont réservés.
        \end{itemize}
      \end{exampleblock}

      \uncover<2->{\par\vspace{0.12cm}\footnotesize Planning statique initial,
        puis dynamique : dates de repli et replanification réactive. Disponible
        sur PC et mobile.}
    \end{column}

    \begin{column}{0.39\textwidth}
      \centering
      \includegraphics[height=4.5cm]{qr_kyst.png}\\[0.25cm]
      {\footnotesize\bfseries Ouvrez KYST}\\[0.1cm]
      {\tiny Données synthétiques, aucune donnée réelle.}
    \end{column}
  \end{columns}
\end{frame}

\begin{frame}{Architecture : un socle standard, hébergé sur place}
  \scriptsize
  \centering
  \begin{tikzpicture}[
      card/.style={draw=McmTerracotta,rounded corners=2pt,fill=white,
                   minimum height=1.10cm,text width=1.85cm,align=center,
                   inner sep=3pt,font=\scriptsize},
      ext/.style={card,draw=McmMuted,dashed},
      flow/.style={{Latex[length=2mm]}-{Latex[length=2mm]},thick,draw=McmTerracotta},
      lab/.style={font=\tiny,text=McmMuted,inner sep=1pt}]
    \node[ext]  (user)  at (0,0)    {\textbf{Navigateur}\\PC ou mobile};
    \node[card] (front) at (3.2,0)  {\textbf{Front}\\SvelteKit};
    \node[card] (api)   at (6.4,0)  {\textbf{API}\\FastAPI};
    \node[card] (db)    at (9.6,0)  {\textbf{Base}\\PostgreSQL};
    \node[ext]  (ai)    at (6.4,-1.75) {\textbf{Modules IA}\\prédiction,\\ordonnanceur};
    \draw[flow] (user)  -- node[lab,above=2pt] {HTTPS} (front);
    \draw[flow] (front) -- node[lab,above=2pt] {interne} (api);
    \draw[flow] (api)   -- node[lab,above=2pt] {SQL} (db);
    \draw[flow] (api)   -- node[lab,right=2pt] {contrats Python} (ai);
    \begin{scope}[on background layer]
      \node[draw=McmSage,dashed,rounded corners=3pt,inner sep=4pt,
            fit=(front)(api)(db)(ai)] (srv) {};
    \end{scope}
    \node[font=\tiny,text=McmSage,anchor=south east] at ([shift={(0,1pt)}]srv.north east)
      {Podman rootless, Ansible};
  \end{tikzpicture}

  \vspace{0.12cm}
  \begin{columns}[T,onlytextwidth]
    \begin{column}{0.49\textwidth}
      \footnotesize Le navigateur ne parle qu'au front ; l'API garde les
      droits et les données.
    \end{column}
    \begin{column}{0.49\textwidth}
      \footnotesize Argon2, cookies HttpOnly/Secure, TLS, conteneurs
      rootless, données pseudonymisées.
    \end{column}
  \end{columns}
\end{frame}

\begin{frame}{Parcours 1/3 : le planning se construit}
  \small
  \centering
  \begin{tikzpicture}[
      c/.style={draw=McmTerracotta,rounded corners=3pt,fill=white,
                text width=2.55cm,minimum height=1.55cm,align=center,
                font=\scriptsize,inner sep=3pt},
      f/.style={-{Latex[length=2.2mm]},thick,draw=McmTerracotta}]
    \node[c] (a) at (0,0) {\textbf{1. Demande}\\patient, spécialité,\\durée estimée};
    \node[c] (b) at (3.25,0) {\textbf{2. Prédiction}\\salle, séjour,\\ambulatoire};
    \node[c] (c) at (6.50,0) {\textbf{3. Dates A/B}\\bloc, lits,\\raisons};
    \node[c] (d) at (9.75,0) {\textbf{4. Confirmation}\\vacation + lit\\réservés};
    \foreach \x/\y in {a/b,b/c,c/d} \draw[f] (\x) -- (\y);
  \end{tikzpicture}

  \par\vspace{0.25cm}\footnotesize Le calcul contraint entre dans le
  parcours clinique, pas dans un tableur séparé.
\end{frame}

\begin{frame}{Parcours 2/3 : le bloc bouge}
  \small
  \renewcommand{\arraystretch}{1.45}
  \setlength{\tabcolsep}{4pt}
  \begin{tabular}{@{}p{3.3cm}p{5.1cm}p{4.9cm}@{}}
    \toprule
    \textbf{Aléa} & \textbf{Ce que fait l'outil} & \textbf{Ce que voit l'équipe} \\
    \midrule
    Salle 1 fermée 10h--12h & replanifie les cas en attente ; les actes
      commencés finissent & le planning avant/après, les cas déplacés \\
    Patient à déplacer & cherche le meilleur créneau compatible & les
      autres patients bougent le moins possible \\
    Chirurgien absent & déplace les actes non commencés & les créneaux
      libérés sont réutilisables \\
    \bottomrule
  \end{tabular}

  \uncover<2->{\par\vspace{0.15cm}\footnotesize Le passé est gelé ; une
    proposition invalide est rejetée et l'ancien plan reste actif.}
\end{frame}

\begin{frame}{Parcours 3/3 : l'urgence arrive}
  \small
  \centering
  \begin{tikzpicture}[x=1.0cm,y=1cm,>=Latex]
    \draw[very thick,draw=McmInk] (0,0) -- (11.0,0);
    \draw[line width=7pt,draw=McmAmber] (3.8,0) -- (5.4,0);
    \foreach \x/\lab in {0/08h,2.2/10h,4.4/12h,7.0/14h,11.0/17h} {
      \draw[draw=McmInk] (\x,-0.12) -- (\x,0.12);
      \node[below=5pt,font=\scriptsize] at (\x,0) {\lab};
    }
    \node[above=16pt,text=McmAmber,font=\scriptsize\bfseries] at (4.6,0)
      {Urgence : fenêtre clinique};
    \draw[line width=3pt,draw=McmSage] (4.0,0.28) -- (5.2,0.28);
    \node[text=McmSage,font=\tiny,anchor=west] at (7.8,0.45)
      {vert : créneau retenu};
  \end{tikzpicture}

  \vspace{0.25cm}
  \begin{columns}[T,onlytextwidth]
    \begin{column}{0.49\textwidth}
      \begin{block}{Ce que fait l'outil}
        \begin{itemize}
          \item insère l'urgence dans sa fenêtre ;
          \item décale le minimum de patients futurs ;
          \item recalcule, valide, puis propose.
        \end{itemize}
      \end{block}
    \end{column}
    \begin{column}{0.49\textwidth}
      \uncover<2->{\footnotesize L'outil ne décide pas seul de décaler un
        patient : la proposition est soumise au praticien, avec son motif.}
    \end{column}
  \end{columns}
\end{frame}

\begin{frame}{Déploiement : sur le serveur de l'établissement}
  \footnotesize
  \begin{columns}[T,onlytextwidth]
    \begin{column}{0.48\textwidth}
      \begin{block}{Ce qui s'installe}
        \begin{itemize}
          \item application web : PC, tablette, téléphone ;
          \item SvelteKit, FastAPI, PostgreSQL ;
          \item conteneurs rootless, déploiement Ansible.
        \end{itemize}
      \end{block}
    \end{column}

    \begin{column}{0.50\textwidth}
      \footnotesize
      Tout reste chez l'établissement.

      \uncover<2->{\par\vspace{0.12cm}Hébergement de données de santé,
        journal d'audit, sauvegardes et droits par rôle.}
    \end{column}
  \end{columns}
  \source{Version actuelle ; ouverture : qualification CE et RGPD.}
\end{frame}

\begin{frame}{Conformité : CE, RGPD et non-substitution}
  \footnotesize
  \begin{columns}[T,onlytextwidth]
    \begin{column}{0.48\textwidth}
      \begin{block}{Marquage CE / MDR}
        \begin{itemize}
          \item Logiciel d'aide à la décision : qualification au titre du
            règlement (UE) 2017/745.
          \item Avant mise sur le marché : marquage CE, dossier technique,
            gestion des risques.
          \item Après mise sur le marché : matériovigilance et mise à jour.
        \end{itemize}
      \end{block}
    \end{column}

    \begin{column}{0.50\textwidth}
      \begin{alertblock}{RGPD et données de santé}
        \begin{itemize}
          \item Minimisation : uniquement les données nécessaires.
          \item Pseudonymisation, hébergement de données de santé (HDS),
            accès par rôle.
          \item Base légale, droits des personnes, journal d'audit.
        \end{itemize}
      \end{alertblock}
    \end{column}
  \end{columns}

  \vspace{0.12cm}
  \begin{columns}[T,onlytextwidth]
    \begin{column}{0.48\textwidth}
      \textbf{Couvert} : dates A/B, réservation, prédictions, replanification.
    \end{column}
    \begin{column}{0.50\textwidth}
      \textbf{À compléter} : audit exhaustif, lissage week-ends et vacances.
    \end{column}
  \end{columns}

  \uncover<2->{\par\vspace{0.18cm}\centering Le praticien confirme chaque
    programme ; la qualification CE et la conformité RGPD restent à instruire.}
  \source{Règlement (UE) 2017/745 ; RGPD (UE) 2016/679.}
\end{frame}
